\section{Extra details to the Outer Gibbs sampler}

This section provides additional technical details on the outer Gibbs sampler described in Section~4.1 \citep{main_article}. 
A summary of the full sampling scheme is given in Algorithm~1. 
We elaborate on two key components: (i) the update of the error variance $\sigma^2$ based on its conjugate inverse-gamma posterior, and (ii) the optional use of the invariant treatment parametrisation introduced by \citet{hahn2020bayesian}, which requires additional Gibbs updates for treatment-specific scaling parameters. 
These details support a more complete understanding and implementation of the full posterior sampling procedure.

The noise is modelled as $\varepsilon_i \sim \mathcal{N}(0, \sigma^2)$, with the error variance $\sigma^2$ assigned a scaled inverse-chi-squared prior:
\begin{equation}
\sigma^2 \sim \text{Inv-}\chi^2(\nu, \psi),
\end{equation}
where $\nu$ denotes the degrees of freedom parameter and $\psi$ the prior scale. We follow \citet{BARToriginal}, setting $\nu = 3$ and choosing $\psi$ so that the prior is centred around a rough estimate of $\sigma^2$.

The variance of the error is sampled from its conditional distribution.
Recall that the prior is $\sigma^2 \sim \text{Inv-}\chi^2(\nu, \psi)$, which corresponds to an inverse-gamma distribution with shape parameter $\frac{1}{2}\nu$ and scale parameter $\frac{1}{2}\nu \psi$. 
By conjugacy, the posterior for $\sigma^2$ is inverse-gamma with updated shape parameter $\frac{\nu + n}{2}$ and scale parameter $\frac{1}{2}\bigl(\nu \psi + \sum_{i=1}^n r_i^2\bigr)$, where $r_i = \log{T^o_i} - \log{\widehat{T}_i}$ denotes the residual for individual $i$, computed using Equation~\eqref{model}.
Detailed derivations and further discussion can be found in \citet[Section 3.2 and 14.2]{Gelman2013BDA}.


We adopt the conventional binary treatment indicator $A$, coded as 1 for treated individuals and 0 for controls, to clearly communicate our model structure and causal assumptions. 
However, the proposed model can also accommodate the invariant parametrisation introduced by \citet{hahn2020bayesian}, which relies on a data-adaptive treatment coding scheme.  
Under the invariant parametrisation, the model is re-expressed as:
\begin{equation}
    \log T(a) = f(x, \hat{e}(x)) + b_a \cdot \tau(x) + \varepsilon,
\end{equation}
where $\varepsilon \sim \mathcal{N}(0, \sigma^2)$ and:
\begin{equation}
    b_1 \sim \mathcal{N}(0, 1/2) \quad \text{and} \quad b_0 \sim \mathcal{N}(0, 1/2).
\end{equation}
In this formulation, the CATE is given by:
\begin{equation}
    \text{CATE}(x) = (b_1 - b_0)\tau(x).
\end{equation}
To implement this approach, two additional update steps are included in the Gibbs sampler to sample $b_1$ and $b_0$. 
Conditional on the prognostic and treatment effect forests, and $\sigma$, these updates correspond to standard linear regression updates. 
Specifically, a linear model where the response is given by $\log T - f(x, \hat{e}(x))$ and the design matrix has two columns (without an intercept), defined as $(A\tau(x), (1-A)\tau(x))$. 
